\begin{abstract}
Tool-using research systems can produce an answer and a detailed execution trace
while leaving unresolved who executed the work, which evidence was available,
whether two agreeing runs were independent, and whether an infrastructure failure
was accidentally admitted as scientific evidence. We define a fail-closed execution
contract that separates request identity, result integrity, executor attribution,
evidence admission and scientific authority. The framework contributes a typed
receipt model, an explicit six-level independence taxonomy and an adversarial
protocol for testing identity substitution, evidence replay, cross-lane leakage,
partial writes and false success. Two contract properties follow from the
definitions: capability failures cannot enter the scientific-evidence set when
admission is type checked, and an agreement label cannot exceed the weakest
verified independence predicate. No protected reliability result is reported.
Unit tests and deterministic source scans, where available, establish only
implementation conformance to their checked cases. The framework therefore makes
execution claims more attributable and falsifiable without treating provenance as
evidence quality or treating agreement as scientific corroboration.
\end{abstract}

\section{Scope and nonclaims}

Computational research increasingly depends on workflows that combine programs,
data services, models, tools and human decisions. Established provenance models
describe the entities, activities and agents involved in producing an object
\citep{moreau2013prov}; workflow profiles and research-object packages make such
records more interoperable and reusable \citep{khan2019cwlprov,soilandreyes2022rocrate,leo2024workflowrun}.
These developments make an execution easier to inspect, but they do not by
themselves answer three distinct questions: whether the recorded executor identity
is authentic, whether an output is scientifically admissible, and whether agreeing
outputs constitute independent corroboration.

This paper addresses that boundary. Its object is a contract between an execution
request, a result receipt and an evidence-admission decision. The contract is
designed for research systems in which a host may call local processes, remote
services, language models or other bounded capabilities. It is not a new workflow
language, provenance ontology, cryptographic signature scheme or research agent.
It instead specifies the minimum predicates that must be checked before a receipt
can support progressively stronger statements about execution and independence.

The paper makes three contributions. First, it separates the identity of a request
from the integrity, attribution and scientific status of its result. Second, it
defines six agreement levels, from deterministic structural agreement to scientific
corroboration, and makes unsupported upward relabeling invalid. Third, it provides
a prospective hostile protocol whose attacks can falsify each guarantee.

The claim boundary is strict. A receipt can show that named bytes were processed
under a declared configuration; it cannot establish that the input evidence was
true or sufficient. A passing unit suite can show implementation conformance; it
cannot estimate reliability outside the checked cases. Two requests can agree while
sharing a model, executor or evidence corpus; agreement alone is not independence.
The framework grants no scientific truth, novelty, publication, deployment or
adoption authority.

